\diarytitle{0117}{極座標変換と球座標変換のヤコビ行列式}

座標変換 $(x,y) \to (u,v)$（あるいは3変数の場合はこれに準ずる）に対するヤコビ行列式は
\[
\frac{\partial(x,y)}{\partial(u,v)} = \det \begin{pmatrix} \dfrac{\partial x}{\partial u} & \dfrac{\partial x}{\partial v} \\[2mm] \dfrac{\partial y}{\partial u} & \dfrac{\partial y}{\partial v} \end{pmatrix}
\]
で定義され、重積分の変数変換公式
\[
\iint_{D} f(x,y)\, dx\, dy = \iint_{D'} f(x(u,v), y(u,v)) \left| \frac{\partial(x,y)}{\partial(u,v)} \right| du\, dv
\]
に用いられる。極座標変換 $x = r\cos\theta,\ y = r\sin\theta$ および球座標変換 $x = r\sin\theta\cos\varphi,\ y = r\sin\theta\sin\varphi,\ z = r\cos\theta$ について、それぞれのヤコビ行列式
$\dfrac{\partial(x,y)}{\partial(r,\theta)}$ および $\dfrac{\partial(x,y,z)}{\partial(r,\theta,\varphi)}$
を、導出過程を示して求めよ。
